\emph{Översikt}
En lista håller reda på värden efter deras \emph{plats}: det första, det
andra, det tredje.  Men det vi oftast vill slå upp något efter är inte en
plats utan ett \emph{namn} --- telefonnumret som hör till Adam, priset
som hör till en artikel, betyget som hör till en student.  Den här
modulen inför \emph{uppslagslistan}
(\foreignlanguage{english}{dictionary}), behållaren som parar ihop en
\emph{nyckel} med ett \emph{värde}.  Vi slår upp, lägger till och ändrar,
hanterar nycklar som saknas, går igenom både nycklar och par, och
jämför uppslagslistan med listan för att kunna välja rätt behållare
utifrån vilka operationer programmet behöver.  Modulen avslutas med en
telefonbok: ett program som byggs upp av användaren, skrivs ut i
bokstavsordning på snygga rader och går att söka i.

\emph{Lärandemål}
Efter denna modul ska studenten kunna:
\begin{restatable}{lo}{DictLOCreate}\label{DictLOCreate}%
  Skapa en uppslagslista, slå upp ett värde med sin nyckel, samt lägga
  till och ändra par --- och förklara varför en nyckel inte är ett
  index.
\end{restatable}% Lo09
\begin{restatable}{lo}{DictLOMissing}\label{DictLOMissing}%
  Hantera att en nyckel saknas, med \mintinline{python}{in},
  med \mintinline{python}{try}/\mintinline{python}{except KeyError}
  eller med \mintinline{python}{get}, och välja mellan sätten.
\end{restatable}% Lo09, Lo04
\begin{restatable}{lo}{DictLOIterate}\label{DictLOIterate}%
  Gå igenom en uppslagslista: dess nycklar, dess värden och dess par
  med \mintinline{python}{items}, samt i sorterad ordning.
\end{restatable}% Lo04
\begin{restatable}{lo}{DictLOChoose}\label{DictLOChoose}%
  Välja mellan lista och uppslagslista utifrån vilka operationer
  programmet behöver, och motivera valet.
\end{restatable}% Lo09
\begin{restatable}{lo}{DictLOProgram}\label{DictLOProgram}%
  Bygga ett program som samlar in data i en uppslagslista, skriver ut
  den läsbart och låter användaren söka i den.
\end{restatable}% Lo02, Lo05, Lo08
\begin{restatable}{lo}{DictLODocs}\label{DictLODocs}%
  Hitta och läsa dokumentationen för uppslagslistor och därifrån
  använda en metod som inte gåtts igenom på föreläsningen.
\end{restatable}% Lo09

\emph{Förkunskaper}
Modulen förutsätter listor och tupler, \mintinline{python}{for}- och
\mintinline{python}{while}-slingor, egna funktioner samt
felhantering med \mintinline{python}{try}/\mintinline{python}{except}.

\emph{Läsning}
Kursens videogenomgång \emph{Behållare: Uppslagslistor} täcker samma
material.
